\defaultfontfeatures{}

\nuaaset{
  thesisid = {032150115},
  title = {舵机测试软件开发},
  author = {易博天},
  college = {自动化学院},
  advisers = {叶永强\quad 教授},
  major = {探测制导与控制技术},
  studentid = {032150115},
  classid = {0322501},
  libraryclassid = {TP273},
  subjectclassid = {081100},
}
\nuaasetEn{
  title = {Literature Review on Servo Test Software Development},
  author = {Yibo Tian},
  college = {College of Automation Engineering},
  majorsubject = {Detection, Guidance and Control Technology},
  advisers = {Prof.~Yongqiang Ye},
  degreefull = {Bachelor of Engineering},
}

\begin{abstract}
本文围绕毕业设计课题“舵机测试软件开发”开展文献综述。面向精密舵机工厂测试这一具体场景，论文重点关注多协议通信、实时数据采集、自动测试组织以及测试结果留档等几类核心需求。在此基础上，本文对执行机构测试平台、工业桌面测试系统架构、RS422/UART 串口可靠通信，以及自动测试与实验设计等方向的代表性研究进行了梳理。综述结果表明，现有工作已经在测试台构建、数据采集与可视化、串口传输机制和局部测试方法等方面积累了较多经验，但面向多型号精密舵机的综合测试软件，仍缺少统一的软件架构和成体系的方法组织。基于这一判断，本文认为后续设计与实现可从分层软件架构、配置化协议解析、预实验驱动的参数生成、测试序列优化以及数据闭环管理等方面展开。上述分析为毕业论文后续研究问题的提出和技术路线的论证提供了文献基础。
\end{abstract}
\keywords{舵机测试软件, RS422通信, 多协议通信, 工业测试系统, 自动测试}

\begin{abstractEn}
This review is written for the graduation project entitled Servo Test Software Development. It examines the needs of precision servo testing in factory settings, with particular attention to multi-protocol communication, real-time data acquisition, automated test execution, and traceable result management. To support the project, representative studies are surveyed from four directions: actuator test benches, industrial desktop test software architectures, reliable RS422/UART serial communication, and automated testing with experimental design. The survey shows that prior work has already provided useful experience in bench construction, data acquisition and visualization, serial transmission, and task-specific testing methods. Even so, a unified software framework for testing multiple servo models is still insufficiently addressed. On this basis, the project is better developed around several connected ideas, namely layered software architecture, configurable protocol parsing, pretest-driven parameter generation, test-sequence optimization, and closed-loop data management.
\end{abstractEn}
\keywordsEn{servo test software, RS422 communication, multi-protocol communication, industrial testing system, automated testing}

\usepackage{subfig}
\usepackage{rotating}
\usepackage[usenames,dvipsnames]{xcolor}
\usepackage{tikz}
\usepackage{pgfplots}
\pgfplotsset{compat=1.16}
\pgfplotsset{table/search path={./fig/}}
\usepackage{ifthen}
\usepackage{longtable}
\usepackage{siunitx}
\usepackage{listings}
\usepackage{multirow}
\usepackage{pifont}
\usepackage{collcell}

\lstdefinestyle{lstStyleBase}{%
  basicstyle=\small\ttfamily,
  aboveskip=\medskipamount,
  belowskip=\medskipamount,
  lineskip=0pt,
  boxpos=c,
  showlines=false,
  extendedchars=true,
  upquote=true,
  tabsize=2,
  showtabs=false,
  showspaces=false,
  showstringspaces=false,
  numbers=left,
  numberstyle=\footnotesize,
  linewidth=\linewidth,
  xleftmargin=\parindent,
  xrightmargin=0pt,
  resetmargins=false,
  breaklines=true,
  breakatwhitespace=false,
  breakindent=0pt,
  breakautoindent=true,
  columns=flexible,
  keepspaces=true,
  framesep=3pt,
  rulesep=2pt,
  framerule=1pt,
  backgroundcolor=\color{gray!5},
  stringstyle=\color{green!40!black!100},
  keywordstyle=\bfseries\color{blue!50!black},
  commentstyle=\slshape\color{black!60}}

\newcommand\cs[1]{\texttt{\textbackslash#1}}
\newcommand\pkg[1]{\texttt{#1}\textsuperscript{PKG}}
\newcommand\env[1]{\texttt{#1}}

\theoremstyle{nuaaplain}
\nuaatheoremchapu{definition}{定义}
\nuaatheoremchapu{assumption}{假设}
\nuaatheoremchap{exercise}{练习}
\nuaatheoremchap{nonsense}{胡诌}
\nuaatheoremg[句]{lines}{句子}
